% !TEX root = ../../Main.tex
\section{Deep Reinforcement Learning}
\label{Section:DeepReinforcementLearning}

Deep Reinforcement Learning (DRL) combines the decision making of reinforcement learning with the function approximation of neural networks. The value functions and policies of \Cref{Section:ReinforcementLearning} are defined over states. In a tabular setting, each state is stored separately, as one entry in a table. This is not possible when the state is a vector of continuous market variables, because there is no limit to the number of distinct states. Replacing the table with a neural network removes that limit. It also lets the agent generalise from states it has seen to states it has not seen.

\subsection{Fundamentals and History}

DRL appeared once it became clear that neural networks could serve as the function approximators that reinforcement learning needed to handle high-dimensional input. The result that started the field was the Deep Q-Network. It learned to play a range of Atari 2600 games directly from raw pixels and reached human-level scores on many of them \cite{mnih_playing_2013, mnih_human-level_2015}. It showed that a single architecture could learn control policies from raw, high-dimensional input, without features designed for each task. This result led to interest in applying the same methods to financial markets. There, the input is also high-dimensional, and the mapping from observation to decision is also unknown. However, an Atari game is stationary and fully observable, and a market is neither. For this reason, results carry over to markets much less directly than the comparison with games suggests.

\subsection{Deep Value-Based RL Algorithms: Deep Q-Networks}

Deep value-based methods use a neural network to approximate the Q-function. This lets them work with state spaces too large to list one by one. Neural Fitted Q Iteration was an early method of this kind. It fitted a network to the Q-values of a stored set of transitions, and fitted it again as more transitions were collected \cite{riedmiller_neural_2005}. The approach became established with the Deep Q-Network (DQN) algorithm \cite{mnih_human-level_2015}, which added two mechanisms that made training stable enough to be practical. The first is a separate target network. This is a delayed copy of the value network, used to compute the learning targets, so that the target does not move as fast as the estimate. The second is experience replay. This is a buffer of past transitions from which mini-batches are drawn at random, which breaks the temporal correlation between consecutive samples. The algorithm used in this work uses both mechanisms.

After DQN, the Double Deep Q-Network (DDQN) dealt with the tendency of DQN to overestimate action values \cite{van_hasselt_deep_2016}. It separates action selection from action evaluation. The online network, the one being trained, chooses the action, and the target network scores it. Deep value-based methods are not used in this work. They select an action by comparing the values of all available actions, and this requires the set of actions to be discrete.

\subsection{Deep Policy-Based RL Algorithms: Deep Policy Gradients}

Deep policy-based algorithms optimise the policy directly, without deriving it from action values. This is what allows them to work with continuous actions. In trading this matters, because a position has a size as well as a direction. A method limited to a discrete set of actions has to decide that size outside the learned policy. The Deep Deterministic Policy Gradient (DDPG) algorithm is a model-free, off-policy, actor-critic method for continuous action spaces \cite{lillicrap_continuous_2015}. It keeps four networks. The actor \( \mu(s | \theta^{\mu}) \) maps a state to a single action. The critic \( Q(s, a | \theta^{Q}) \) estimates the value of a state-action pair. The other two are delayed copies of these, the target actor \( \mu' \) and the target critic \( Q' \). Transitions are stored in a replay buffer, and mini-batches of \( N \) transitions are sampled from it at random. The critic is trained in the same way as in DQN. It minimises the squared error between its estimate and a target built from the observed reward and the value of the next state \cite{lillicrap_continuous_2015}:

\begin{equation}
y_i = r_i + \gamma \, Q'\!\left(s_{i+1}, \mu'(s_{i+1} | \theta^{\mu'}) \,\middle|\, \theta^{Q'}\right)
\label{Equation:DDPGTarget}
\end{equation}

\begin{equation}
L = \frac{1}{N} \sum_{i} \left( y_i - Q(s_i, a_i | \theta^{Q}) \right)^2
\label{Equation:DDPGCriticLoss}
\end{equation}

The actor is updated with the deterministic policy gradient, which follows the direction in which the value estimated by the critic increases. Because the action is continuous, the gradient of the value with respect to the action is defined, and the chain rule carries it back to the parameters of the actor \cite{lillicrap_continuous_2015}:

\begin{equation}
\nabla_{\theta^{\mu}} J \approx \frac{1}{N} \sum_{i} \nabla_{a} Q(s, a | \theta^{Q}) \big|_{s = s_i,\, a = \mu(s_i)} \; \nabla_{\theta^{\mu}} \mu(s | \theta^{\mu}) \big|_{s_i}
\label{Equation:DeterministicPolicyGradient}
\end{equation}

The target networks are not simply copied from the trained networks. After every update, they are moved a small fraction \( \tau \ll 1 \) towards them. This keeps the learning targets almost stationary, but it also slows down how fast they follow improvements \cite{lillicrap_continuous_2015}:

\begin{equation}
\theta' \leftarrow \tau \theta + (1 - \tau) \theta'
\label{Equation:SoftUpdate}
\end{equation}

The policy is deterministic, so it does not explore by itself, and the exploration problem raised in \Cref{Section:ReinforcementLearning} comes back here. DDPG solves it by adding a noise process \( \mathcal{N} \) to the action during training, so the action actually executed is \( \mu(s | \theta^{\mu}) + \mathcal{N} \). The original work uses an Ornstein-Uhlenbeck process. This process gives noise that is correlated in time, instead of independent at each step. The reason given is that a physical system with momentum explores better when a disturbance lasts for a while than when it changes sign at every step \cite{lillicrap_continuous_2015}. \Cref{Codes:DeepDeterministicPolicyGradient} shows the complete procedure.

% !TEX root = ../Main.tex
\begin{algorithm}[htb!]
\DontPrintSemicolon
\KwRequire{an actor $\mu(s | \theta^{\mu})$, a critic $Q(s, a | \theta^{Q})$, a discount factor $\gamma$, a soft update rate $\tau \ll 1$ and a minibatch size $N$}
Randomly initialise the critic parameters $\theta^{Q}$ and the actor parameters $\theta^{\mu}$\;
Initialise the target parameters $\theta^{Q'} \leftarrow \theta^{Q}$ and $\theta^{\mu'} \leftarrow \theta^{\mu}$, and the replay buffer $\mathcal{B}$\;
\For{episode $= 1, M$}{
    Initialise a random process $\mathcal{N}$ for action exploration\;
    Receive the initial observation state $s_1$\;
    \For{$t = 1, T$}{
        Select the action $a_t = \mu(s_t | \theta^{\mu}) + \mathcal{N}_t$ from the current policy and the exploration noise\;
        Execute the action $a_t$ and observe the reward $r_t$ and the new state $s_{t+1}$\;
        Store the transition $(s_t, a_t, r_t, s_{t+1})$ in $\mathcal{B}$\;
        Sample a random minibatch of $N$ transitions $(s_i, a_i, r_i, s_{i+1})$ from $\mathcal{B}$\;
        Set $y_i = r_i + \gamma Q'(s_{i+1}, \mu'(s_{i+1} | \theta^{\mu'}) | \theta^{Q'})$\;
        Update the critic by minimising the loss: $L = \frac{1}{N}\sum_i (y_i - Q(s_i, a_i | \theta^{Q}))^2$\;
        Update the actor policy using the sampled policy gradient:\;
        $\nabla_{\theta^{\mu}} J \approx \frac{1}{N}\sum_i \nabla_a Q(s, a | \theta^{Q})|_{s = s_i, a = \mu(s_i)} \nabla_{\theta^{\mu}} \mu(s | \theta^{\mu})|_{s_i}$\;
        Update the target networks:\;
        $\theta^{Q'} \leftarrow \tau \theta^{Q} + (1 - \tau) \theta^{Q'}$\;
        $\theta^{\mu'} \leftarrow \tau \theta^{\mu} + (1 - \tau) \theta^{\mu'}$\;
    }
}
\caption{Deep Deterministic Policy Gradient \cite{lillicrap_continuous_2015}.}
\label{Codes:DeepDeterministicPolicyGradient}
\end{algorithm}

DDPG is powerful, but it is not reliable. It has no guarantee of convergence except under restrictive conditions. It is also known to be sensitive to its hyperparameters and to the random seed, so two runs of the same configuration can end with different policies. A market makes this worse. It is non-stationary and only partially observable, which means the agent does not see the full state, and the guarantees that apply to finite Markov Decision Processes do not hold. These limits belong to the method itself. They are not problems that can be fixed. \Cref{Chapter:ExperimentalResults} therefore reports how the results vary, instead of reporting a single run.